\documentclass[letterpaper]{article}
\PassOptionsToPackage{unicode=true}{hyperref} % options for packages loaded elsewhere
\PassOptionsToPackage{hyphens}{url}
\usepackage{graphicx} % Required for inserting images


\def\tightlist{}


\graphicspath{ {../images/} }
\usepackage[margin=1in]{geometry}
\usepackage{tikz}
\usepackage[dvipsnames]{xcolor}
\usetikzlibrary{calc}




\usepackage[utf8]{inputenc}
\usepackage[T1]{fontenc}
\usepackage{lmodern}
\usepackage[english]{babel}

\usepackage{longtable}
\usepackage{booktabs}

% Useful packages
\usepackage{hyperref}
\usepackage{amsmath, amssymb}
\usepackage{fancyhdr}
\usepackage{setspace}
\usepackage{tocloft}
\usepackage{titlesec}


\usepackage{unicode-math}



% Code highlighting (Pandoc uses listings or minted, but listings is safer without shell escape)
\usepackage{listings}
\lstset{
  basicstyle=\ttfamily\small,
  breaklines=true,
  frame=single,
  backgroundcolor=\color[gray]{0.95}
}

\usepackage{setspace}


%Modifiable Commands; defaults set to none
\newcommand{\titleDOC}{}
\newcommand{\authorDOC}{Daniel Sabalakov}
\newcommand{\dateDOC}{\today}
\newcommand{\classDOC}{}
\newcommand{\emailDOC}{danielsabalakov@gmail.com}
% DEFAULT QUOTE IS BY ISSAC NEWTON. IF YOU DON'T WANT IT, DO: quote: @none
\newcommand{\quoteDOC}{``\textit{Truth is ever to be found in the simplicity, and not in the multiplicity and confusion of things}'' - Issac Newton}

% Header/Footer
\pagestyle{fancy}
\fancyhf{}
\rhead{\thepage}
\lhead{\leftmark}
\fancyhead[R]{\thepage} % Place page number in top right corner




% %Title Arguments
% \title{\TITLEOFDOCUMENT}
% \author{\AUTHOROFDOCUMENT}
% \date{\today}


%
% ANYTHING BELOW HERE CAN GET PROCEDURALLY GENERATED
%
% BEGINNING OF DOCUMENT
%
%
%
\begin{document}

%titlepage








\renewcommand{\titleDOC}{3.2 Eliud Kipchoge}

\renewcommand{\quoteDOC}{``\textit{May the Force Be With You}'' \textemdash{} General Jan Dodonna}

\renewcommand{\classDOC}{PE Running}

\renewcommand{\authorDOC}{Daniel Sabalakov}





\begin{titlepage}
  \titleDOC
  \classDOC
  \dateDOC
  \authorDOC
\end{titlepage}

\section{Eliud Kipchoge}\label{eliud-kipchoge}

 \begin{doublespace} 

Eliud Kipchoge, arguably the best runner of all time, is a 41 year old distance runner from Kenya. He is 5\textquotesingle7 (1.7 meters) tall, and he weighs 115 pounds (52 kg). He was born in Kenya and competed in many Olympic marathons. He was the world record holder in the marathon from the years of 2018-2023. He has ran four out of the ten fastest marathons in history. In 2003, he won the "junior race at the World Cross Coutnry Championships and setting a world junior record for the 5000 m." (Wikipedia). He then became the world champion at the 2003 World Championships, quickly followed by two bronze metals in the 2004 Olympic Games and the World Indoor Championships. However, he is not primarily known for his 5k times.

In 2012, Eliud switched to road running. He ran the half marathon in 59:25, which was the second fastest debut for the half marathon ever (Wikipedia). He won the 2013 Hamburg Marathon, and then he won the Chicago Marathon. He then proceeded to become a World Marathon Major 5 more times: in 2016, 2017, 2018, 2019, and 2022. In the current day, he still runs marathons. He placed seventeenth in the 2025 New York City Marathon, completing all six World Marathon Major courses. Eliud now plans to "run seven marathons across seven continents over a period of two years to raise awareness about sustainability and education." (Wikipedia)

Kipchoge trains at 80\% of his effort on Tuesday, Thursday, and Saturday. The rest of the week, he runs at 50\%. (Run) Eliud also perform 60 minutes of strength training, "Twice a week, Kipchoge and his training partners perform a 60-minute session of strength and mobility exercises using yoga mats and resistance bands. The exercise program focuses on the posterior chain, particularly the glutes, hamstrings, and core muscles." (Run)

Eliud Kipchoge is the epitome of human athletisism. From winning multiple World Championships in the 5k, to then becoming an elite distance Marathoner, Eliud is undoubtedly one of the best distance runners of all time.

 \end{doublespace} 

\section{Sources}\label{sources}

\begin{itemize}
\tightlist
\item
  https://en.wikipedia.org/wiki/Eliud\_Kipchoge
\item
  https://run.outsideonline.com/training/training-plans/marathon/eliud-kipchoge-marathon-workout-training-principles/?scope=anon
\end{itemize}



\end{document}
